% Generated draft: requires mathematical and source review.
\documentclass[11pt]{article}
\usepackage[margin=25mm]{geometry}
\usepackage{fontspec}
\setmainfont{Latin Modern Roman}
\usepackage{amsmath,amssymb,mathtools,unicode-math}
\setmathfont{Latin Modern Math}
\usepackage{enumitem,needspace,xurl,hyperref}
\setlength{\parindent}{0pt}
\setlength{\parskip}{5pt}
\setlength{\emergencystretch}{3em}
\newcommand{\N}{\mathbb{N}}
\newcommand{\Z}{\mathbb{Z}}
\newcommand{\Q}{\mathbb{Q}}
\newcommand{\R}{\mathbb{R}}
\newcommand{\C}{\mathbb{C}}
\newcommand{\F}{\mathbb{F}}
\newcommand{\A}{\mathbb{A}}
\newcommand{\PP}{\mathbb{P}}
\newcommand{\E}{\mathbb{E}}
\newcommand{\eps}{\varepsilon}
\newcommand{\dd}{\,\mathrm{d}}
\begin{document}
\section*{UnsolvedMath ID: 30004255}
\textbf{Source problem code:} catalogue entry
\subsection{Definitions and mathematical statement}
Write \(X\) for the underlying curve implicit in the source's geometric Langlands setting, and \(G\) for its group. The source does not specify either object or the base field. Let \(\operatorname{Bun}_G(X)\) denote the moduli stack of principal \(G\)-bundles on \(X\), and write \(\operatorname{LocSys}_{\check G}(X)\) for the moduli object of local systems on \(X\) for the intended Langlands dual group \(\check G\). The meaning of local system and the precise moduli object depend on a geometric Langlands setting that the record leaves unspecified.

For a quantization \(M\) as in the source, denote its geometric theta series by \(\Theta_M\); the hypothesis is that \(\Theta_M\) is a distinguished \(D\)-module on \(\operatorname{Bun}_G(X)\). The terms ``quantization,'' ``geometric theta series,'' and ``distinguished'' are not defined in the supplied record. The notation \(\check M\) denotes the intended dual object appearing in the source; its construction and its relationship to \(M\) are also unspecified.

The intended question is whether, for every such \(M\), its geometric Langlands dual is a sheaf on \(\operatorname{LocSys}_{\check G}(X)\) arising from a suitable coherent quantization of \(\check M\). Schematically, writing \(\mathsf L\) for the intended geometric Langlands correspondence or transform, the proposed relationship is
\[
\mathsf L(\Theta_M)\simeq\mathcal F_{\check M},
\qquad
\mathcal F_{\check M}\text{ arises from a suitable coherent quantization of }\check M.
\]
This display records the proposed relationship, not a defined functor or a constructed sheaf. A precise quantified conjecture requires the admissible class of \(M\), the geometric Langlands setting and categories, the meaning of \(\mathsf L\), the construction of \(\check M\), and the procedure producing \(\mathcal F_{\check M}\) from a coherent quantization. In particular, the source does not specify whether ``sheaf'' means a quasi-coherent, coherent, ind-coherent, or another kind of object. These unresolved qualifications prevent a fully determined mathematical statement.
\subsection{Short English statement}
The question asks whether the geometric Langlands dual of a distinguished geometric theta \(D\)-module associated with a quantization \(M\) is a sheaf produced from a suitable coherent quantization of the intended dual object \(\check M\). The geometric setting and the constructions needed to make this precise are missing from the supplied record.
\subsection{Sources}
\begin{itemize}
\item Supplied local UnsolvedMath registry record, permanent ID \(30004255\), title ``Geometric Langlands Duals of Quantized Theta Sheaves,'' source problem code \texttt{catalogue entry}, in the Oberwolfach Reports Open Problems collection. The record gives proposed year 2019 and identifies no individual proposer.
\item Supplied source citation: \emph{New Developments in Representation Theory of p-adic Groups} (2020), Oberwolfach Reports, pp.~2739--2819, DOI \texttt{10.4171/owr/2019/44}. \url{https://doi.org/10.4171/owr/2019/44}. The supplied background describes Oberwolfach Reports as reports of mathematical workshops and their open-problem discussions. No problem-specific page or passage is supplied.
\item Supplied literature source, labeled ``OWR 2019 44'': \url{https://ems.press/content/serial-article-files/46824}.
\item Supplied literature source, labeled ``Feng--Wang 2025'': \url{https://link.springer.com/article/10.1007/s00039-025-00702-4}. The record supplies no article title or problem-specific result from this reference.
\item Literature assessment dated exactly \texttt{2026-08-22}: the supplied curation labels the problem ``partially solved'' and describes the global theta-sheaf statement as part of the Ben-Zvi--Sakellaridis--Venkatesh conjectural program, with semiclassical and selected low-rank examples providing evidence rather than a general proof. This is supplied dated curation derived from the OWR record, not an independently verified literature assessment; its status and cited claims require independent verification.
\end{itemize}
\end{document}
